\author{OTS}
\documentclass[14pt]{../../inc/mybib}

\setincpath{../../inc/}

\usepackage{bible_style}
\usepackage{textmakro}
\graphicspath{{../../assets/images/}}
\usepackage{header}

\newcommand{\Name}{Nathanael}

% ensure scrlayer-scrpage has sufficient footheight
\setlength{\footheight}{10.4pt}
\setlength{\headheight}{0.4pt}

\begin{document}
\begin{bibelbox}{SCHL}{Kol}{3:15-17}
    Lasst das Wort des Christus reichlich in euch wohnen in aller Weisheit; lehrt und ermahnt einander und singt mit Psalmen und Lobgesängen und geistlichen Liedern dem Herrn lieblich in eurem Herzen. Und was immer ihr tut in Wort und Werk, das tut alles im Namen des Herrn Jesus und dankt Gott, dem Vater, durch ihn.
\end{bibelbox}
\section{Begrüssung}
Ich möchte auch euch alle recht herzlich zum Gottesdienst begrüssen. Schön, dass ihr gekommen seid, um unserem \herr N zu danken, ihn zu loben und zu preisen.

Auch dich \Name{} will ich herzlich begrüssen. Schön, dass du wieder bei uns bist und mit uns Gottesdienst feierst.
\begin{itemize}
    \item[] \beten{}
    \item[] Und anschliessend singen wir zusammen das Lied
    \item[] \lied{Kommt, stimmt doch mit uns ein}.
\end{itemize}

\section{Informationen}
\begin{itemize}
    \item \bt[infos]{Bibel und Gebetsabend:} Do, 16.April.2026 20:00 Uhr Bibel und Gebetsabend mit Samuel Rindlisbacher Markus 6,30-44
    \item \bt[infos]{Nächster Gottesdienst:} So, 19.April.2026 14:45 Uhr mit Fredy Peter    
    \item \bt[infos]{Vortrag:} Morgen Abend wird uns Nathanael Winkler im Rahmen seiner Schweizer Tournee hier in diesen Räumen einen Vortrag über Gog und Magog: \flqq{} Alte Prophezeiung, neue Geopolitik\frqq{} halten. Ladet Freunde und Kollegen ein.
    \item \bt[infos]{Kollekte:} Die Kollekte wird hinten in der grünen Box gesammelt und wird vollumfänglich für den Bau dieser Gemeinde eingesetzt.
\end{itemize}

\section{ Input }
\begin{spacing}{1.5}
    \begin{block}[Gott hat die Kontrolle]
        Das war das Motto dieser Osterkonferenz. Ich vermute mal, dass ein Teil von hier die Vorträge gehört haben. Familie Todt und Familie Schmid war vor Ort und für waren diese Vorträge sehr ermutigend. Vor allem für die Arbeit hier vor Ort. Aber nicht nur dafür sondern auch für die Zukunft. Gott hat die Kontrolle. Wir sind ja schon seid einiger Zeit in der Heilsgeschichte Gotte unterwegs und wenn man die Bibel aus dieser Perspektive liest, dann sieht man sieht man Gottes wirken, seid Beginn der Schöpfung.

        Letztes Mal haben wir uns den grossen weissen Thron in Offenbarung 20 angesehen. Wir haben gesehen, dass Jesus Gericht über alle jemals gelebten Menschen gehalten hat. Und jetzt? Ist das alles gewesen? Nein, es geht weiter. Das ist nicht das Ende. Jetzt  beginnt es erst.

        Manchmal frage ich mich, wieso Gott diesen Aufwand mit uns Menschen überhaupt betreibt. Ich glaube, dass Gott die Ewigkeit nur mit Menschen verbringen will, die ihn von Herzen lieben. Seid dem Sündenfall, hat Gott Menschen ausgesondert, die freiwillig ihr Herz für Gott geöffnet haben.

        Um mit diesen Menschen zusammen die Ewigkeit verbringen zu können, schafft Gott alles Neu. Und wenn ich \betonung{NEU} sage, dann meine ich wirklich \betonung{NEU}. Es wird eine neue Erde geben, einen neuen Himmel und ein neues Jerusalem. Johannes beschreibt wie diese neue Welt aussehen wird, und es klingt einfach fantastisch. Es wird keinen Tod mehr geben, keine Tränen, keine Schmerzen, keine Krankheiten, keinen Hunger, keinen Krieg, keine Sünde mehr. Es wird eine perfekte Welt sein. Und das alles in der Gegenwart Gottes. 

        Um uns nochmals aufzurütteln schreibt Johannes in 21.8 über die Feiglinge, Ungläubigen, Abtrünnigen, Mörder, Sexsünder, Zauberer, Götzendiener und Lügner, dass sie alle im Feuersee landen ihren Anteil am zweiten Tod haben werden. Es wird also zwei bereiche geben, die neue Welt und der Feuersee.
        \bibleverse{Offb}(20:1-10) 
        Thomas hat uns hier zwei Sonntage das Tausendjährige Reich nahegebracht. Darum brauche ich euch heute nicht damit zu langweilen, sondern will dies zum Anlass nehmen, etwas über Irrlehren zu sagen. Wie auch schon Thomas gesagt hat, glaube auch ich, dass diese 1000 Jahre wörtlich zu nehmen sind. Wieso sollte diese Zahl sonst so oft erwähnt werden? Manche Ausleger sagen, dass die Zahl 1000 \betonung{Vollkommenheit} bedeutet. Diese Sichtweise wird Amillennialismus genannt, also \betonung{kein Millennium}. Eine Theorie ist auch, dass wir jetzt schon im Millennium sind. Die kath. Kirche vertritt diese Ansicht. Sie sagt, Jesus kommt wieder, wenn hier auf der Erde Friede und Ruhe herrscht.

        Andere Theorien sind, dass das tausendjährige Reich mit dem ersten Kommen Jesu begonnen hat. Das Königreich ist schon errichtet und jeder, der Christus von Herzen aufnimmt, trägt das Königreich in sich. Der Teufel wurde mit dem Tod und der Auferstehung Christi gebunden. Es sind ziemlich bekannte YouTube-Prediger, die diese Ansicht vertreten. So z\.B auch Olaf Latzel, Tobias Riemenscheider, Willy Zorn, Jürgen Fischer.
        Ich habe auf YouTube eine Debatte gehört, von einem Millennialisten und einem Amillennialisten. Beides bibeltreue Christen, und beide haben mit der Bibel argumentiert, aber als Amillennialist muss man viele Bibelstellen vergeistlichen. Für mich als Millennialist ist darum eine buchstäbliche Auslegung bei Weitem einfacher. Wir sollten aber immer beachten, dass jemand, nur weil er zu diesem Thema eine andere Sicht hat, darum noch kein Irrlehrer ist. Irrlehrer sind solche, die das falsche Evangelium verkünden.\\  
        \textbf{Biblische Merkmale für Irrlehrer}
        \begin{enumerate}
            \item \textbf{Falsches Evangelium}: Wer ein anderes Evangelium bringt als das apostolische, steht unter scharfer Warnung (\bibleverse{Gal}(1:6-9)).
            \item \textbf{Falsches Christusbild}: Wer Jesus nicht gemäß der biblischen Wahrheit bekennt (die Menschwerdung, seine Herrschaft, sein Werk) (\bibleverse{IJoh}(4:1-3), \bibleverse{IIJoh}(7:)).
            \item \textbf{Verdrehung der Gnade:} Gnade wird als Freibrief zur Sünde missbraucht (\bibleverse{Jud}(4:)).
            \item \textbf{Leugnung zentraler Glaubensinhalte:} z. B. Auferstehung oder Herrschaft Christi (\bibleverse{IITim}(2:17-18), \bibleverse{IIPetr}(2:1)).
            \item \textbf{Schlechter Fruchtcharakter:} Habgier, Ausbeutung, Stolz, Unmoral, Spaltung. 
            (\bibleverse{Matt}(7:15-20), \bibleverse{IIPetr}(2:2-3), \bibleverse{Titus}(3:10)). Der Lebensstil eines Predigers kann viel über ihn aussagen.
        \end{enumerate}        
        Bevor Jesus in Macht und Herrlichkeit wiederkommt, werden viele solcher Irrlehrer auftreten. In der heutigen Zeit mit dieser weltweiten Vernetzung ist es einfach, ein falsches Evangelium schnell zu verbreiten, und es ist leider nicht immer so offensichtlich.

        Ja, Gott liebt alle Menschen, er liebt dich und mich. Gott hat aber auch Gefühle, er wird zornig, kann enttäuscht werden, ist traurig, und das nur, weil Menschen ihn dauernd enttäuschen. Gott liebt uns so sehr, dass er seinen Sohn für uns am Kreuz geopfert hat. \betonung{Wie danken wir ihm dafür?} Er verlangt ja nur, dass wir ihn lieben, mehr nicht. Lieben wir Gott? Oder hoffen wir nur, nicht in die Hölle zu kommen?

        Und für alle, die jetzt denken, dass der Mensch doch im Grunde seines Herzens gut ist, sehen wir hier im Millennium den Beweis, dass dem nicht so ist. Zu der Zeit ist Satan im Abgrund gebunden. Er hat keine Macht. Jesus ist der Herrscher und herrscht gerecht. Es wird Regeln und Gesetze geben, die absolut gerecht gegen jeden sind. Wir lesen, dass sich Menschen gegen diese Regeln auflehnen werden. Wer z\.B mit 100 Jahren stirbt, der hat gesündigt. Andere werden murren. Wiersbe schreibt dazu:
        \begin{block}[Zitat von Warren W. Wiersbe:]
            In gewissem Sinne wird das tausendjährige Reich all das \betonung{zusammenfassen}, was Gott über das Herz des Menschen in den verschiedenen Geschichtsabschnitten gesagt hat. Es wird eine Herrschaft des Gesetzes (das gerecht sein wird) sein, und doch wird das Gesetz das sündige Herz des Menschen nicht verändern können. Der Mensch wird sich immer noch gegen Gott (Jesus) auflehnen.
        \end{block}        
        Obwohl es keine Krankheiten, keinen Hunger und keine Kriege mehr gibt, wenden sich Menschen von Jesus ab. Auch perfekte Bedingungen bringen also keine guten Herzen hervor. Darum wird Gott zum Abschluss Satan nochmals auf die Menschen loslassen. Der Satan wird diese Abtrünnigen um sich scharen, um gegen Jesus anzutreten. Gott wird dieses Mal kurzen Prozess machen, Off 20:7-10. Das Böse wird endgültig vernichtet. Was danach passiert, hört ihr dann am nächsten Sonntag.   
        
    \end{block}
    \lied{Glauben heisst vertrauen}.
\end{spacing}
\section{Predigt}
\textbf{Nach der Predigt}\newline
Danken für die Predigt.
% \section{Abendmahl}
% \begin{itemize}
%     \item [] Beten für das Brot
%     \item [] Beten für den Wein
% \end{itemize}
\section{Abschluss}
\begin{itemize}
    \item [] \lied{Herr du gibst Hoffnung}.
    \item [] \beten{}
\end{itemize}
\begin{bibelbox}{SCHL}{IIKor}{13:13}
    Die Gnade des Herrn Jesus Christus und die Liebe Gottes und die Gemeinschaft des Heiligen Geistes sei mit euch allen! Amen
\end{bibelbox}
\end{document}